\documentclass[11pt,a4paper]{article}
\usepackage[margin=27mm]{geometry}
\usepackage{amsmath,amssymb,amsthm,hyperref}
\newtheorem{theorem}{Theorem}
\newtheorem{lemma}[theorem]{Lemma}
\newtheorem{corollary}[theorem]{Corollary}
\newcommand{\gen}[1]{\langle #1\rangle}
\title{Translation kernels in affine reflection groups}
\author{Human Compression project}
\date{}
\begin{document}\maketitle

\begin{abstract}
A normal subgroup of an irreducible affine Coxeter group which meets
every maximal standard parabolic trivially consists of translations.
The proof uses the commutator of an isometry with a wall reflection:
nonparallel mirrors would give a forbidden finite-order element.
This supplies Corollary~5.1 of the fixed hypertope manuscript under its
Section~5 restriction to universal spherical residues. We separate that
hypothesis from the weaker assertion that the residue diagrams are finite.
\end{abstract}

\section{The geometric input}

Let $W=\gen{s_i:i\in I}$ be an irreducible affine Coxeter group in its
faithful action by Euclidean isometries on $E=\mathbb R^d$, where
$d\geq1$ and $|I|=d+1$. Write $M_i$ for the mirror of $s_i$ and
$W_i=\gen{s_j:j\ne i}$. We retain the usual affine Coxeter realization:
$W$ is discrete, its closed fundamental alcove is a bounded simplex,
every point is carried into that alcove by $W$, and the stabilizer of a
point in the alcove is generated by the reflections in its containing
walls. These are standard geometric Coxeter facts, not a classification
of its normal subgroups. The linear normals to the walls span $E$;
otherwise the alcove would contain a line. Their nonorthogonality graph
is connected, as is the irreducible affine Coxeter diagram.

The translation subgroup is
\[
 T=\{w\in W:w(x)=x+b\text{ for some }b\in E\}.
\]
No finite-index assumption will be placed on a subgroup of $W$.

\begin{lemma}\label{torsion}
If $N\triangleleft W$ and $N\cap W_i=1$ for every $i\in I$, then
$N$ is torsion-free.
\end{lemma}
\begin{proof}
Suppose $g\in N$ has finite order $m$. For any $x\in E$, the average
$m^{-1}\sum_{k=0}^{m-1}g^k x$ is fixed by $g$, since isometries are
affine maps. Carry this point into the closed fundamental alcove.
The conjugate of $g$ lies in $N$ by normality and in the stabilizer
of that point. The containing walls are a proper subset of all alcove
walls: a nondegenerate simplex has no point on every facet.
Their generated group is therefore contained in some $W_i$.
The assumed intersection forces the conjugate, and hence $g$, to be
the identity.
\end{proof}

We also record why discreteness pays the other direction used below.
The stabilizer of a point in a discrete isometry group is finite:
after moving the point to the origin it is a discrete subgroup of the
compact orthogonal group. If it were infinite, a convergent subsequence
would give distinct group elements whose quotients tend to the identity,
contradicting discreteness. Thus every element of $W$ fixing a point
has finite order.

\section{A reflection commutator}

\begin{theorem}\label{translation}
Let $W$ be as above. Every torsion-free normal subgroup of $W$ is
contained in $T$. In particular,
\[
 N\triangleleft W,\qquad N\cap W_i=1\ (i\in I)
 \quad\Longrightarrow\quad N\leq T.
\]
\end{theorem}
\begin{proof}
Let $g\in N$ and write $g(x)=Ax+b$ with $A$ orthogonal.
For each $i$, normality gives
\[
                c_i=g s_i g^{-1}s_i\in N.
\]
Indeed $s_i g^{-1}s_i\in N$. The conjugate $g s_i g^{-1}$ is
the reflection in $gM_i$.

Assume first $d\geq2$. If $gM_i$ and $M_i$ are not parallel, they
intersect. The product of their reflections fixes every point of that
intersection. It therefore has finite order in the discrete group $W$.
The mirrors are distinct, so their reflections are distinct and their
product is not the identity. This contradicts torsion-freeness of $N$.
Consequently $gM_i$ is parallel to $M_i$ for every $i$.

Choose a nonzero normal $u_i$ to each $M_i$. Parallelism and
orthogonality of $A$ imply $Au_i=\varepsilon_i u_i$ with
$\varepsilon_i\in\{1,-1\}$. If $\langle u_i,u_j\rangle\ne0$, then
\[
 \langle u_i,u_j\rangle
 =\langle Au_i,Au_j\rangle
 =\varepsilon_i\varepsilon_j\langle u_i,u_j\rangle,
\]
so $\varepsilon_i=\varepsilon_j$. Connectivity makes all signs equal.
Since the normals span $E$, this proves $A=I$ or $A=-I$.
In dimension one these are already the only orthogonal linear maps,
so the same conclusion holds without the nonparallel-mirror argument.

If $A=-I$, then $g(x)=-x+b$ is a nonidentity involution, fixing $b/2$.
This is again impossible in $N$. Hence $A=I$, and $g$ is a translation.
The final implication follows from Lemma~\ref{torsion}.
\end{proof}

The theorem uses normality twice, and irreducibility exactly in passing
from preservation of the wall-normal lines to a common sign. It does not
assert the conclusion for a general torsion-free subgroup, or for a
reducible Euclidean reflection group.

\section{The hypertope application and its precise domain}

Let $\mathcal H$ be a regular hypertope with the irreducible affine
Coxeter diagram of $W$. Its distinguished chamber involutions give a
marked epimorphism
\[
                 \pi:W\longrightarrow G=\operatorname{Aut}_I(\mathcal H).
\]
Here we use the acquired residue-stabilizer proof: the involutions
generate $G$, have the exact diagram labels, and each maximal marked
parabolic $G_i$ acts faithfully and regularly on the corresponding
residue. These conclusions do not assume that $\pi$ is locally injective.

\begin{corollary}\label{hypertope}
Suppose every maximal residue of $\mathcal H$ is the universal spherical
Coxeter complex with its inherited diagram. Then $N=\ker\pi$ lies in
$T$, and $\mathcal H$ is the incidence quotient of the universal affine
Coxeter complex by $N$.
\end{corollary}
\begin{proof}
The restriction $W_i\to G_i$ is onto. The universal spherical residue
hypothesis identifies its finite domain and image orders: both are the
number of chambers of that spherical complex. Hence the restriction is
injective, so $N\cap W_i=1$. Theorem~\ref{translation} gives $N\leq T$.

It remains to identify the incidence quotient, rather than only its
group. An object of type $i$ in the universal complex is a left coset
$wW_i$. Its $N$-orbit corresponds exactly to the coset
$\pi(w)G_i$. Two quotient objects are incident when some representatives
of their orbits are incident in the universal complex. This is equivalent
to the two corresponding cosets in $G$ having a common representative:
a common image lifts to $W$, and the discrepancy is absorbed by $N$.
Thus the incidence quotient is the coset incidence system of $G$ with
the subgroups $G_i$. The acquired reconstruction theorem for the actual
regular hypertope identifies it with $\mathcal H$.
\end{proof}

This is a necessary-kernel statement for an already existing hypertope.
An arbitrary normal translation subgroup need not give a hypertope.
In particular, translations automatically avoid the finite parabolics,
but can meet their products. The exact product-kernel condition and
flag transitivity still have to be checked for a proposed quotient, as
shown in the separately acquired correction to source Theorem~3.2.
No converse existence or full period-lattice classification is asserted.

There is a genuine terminology issue in the fixed source. On page~5,
``spherical'' is defined by the finite Coxeter diagram, which also
allows projective quotients. On page~10, Section~5 explicitly distinguishes
spherical residues from projective residues and restricts attention to
the former. Corollary~\ref{hypertope} proves source Corollary~5.1 in
this latter, locally universal sense. Finite residue diagrams alone do
not imply the injectivity $W_i\to G_i$ used above; a projective residue
has only half as many chambers. The source itself mentions a Euclidean
example with projective residues on page~9. We neither silently exclude
those examples from the earlier definition nor claim their kernels have
been classified by the present argument. The same local hypothesis must
be supplied when using the injectivity assertion in source Theorem~3.4.

\paragraph{Attribution and verification.}
The translation-kernel conclusion is the source's result, attributed
there to Maxwell's affine normal-subgroup theorem. We have not acquired
the full primary proof of that theorem and do not use it here. The new
argument is the short reflection-commutator proof above, together with
the exact residue interface. The affine realization and alcove-stabilizer
facts remain named standard inputs. The supplementary exact matrix checks
exercise the commutator step in two plane reflection groups; the uniform
theorem is proved by the written argument, not by those finite examples.
No historical priority, open-conjecture, new Lean or measurement claim is made.

\begin{thebibliography}{9}
\bibitem{source} M. E. Fernandes, D. Leemans and A. Ivi\'c Weiss,
\emph{An exploration of locally spherical regular hypertopes}.
Fixed sixteen-page revised manuscript, catalogue L033, pages 5, 7, 9--11.
\bibitem{maxwell} G. Maxwell, \emph{The normal subgroups of finite and affine
Coxeter groups}, Proc. London Math. Soc. 76 (1998), 359--382.
The source attributes its affine conclusion to Theorem 0.2; this theorem
is not a premise of the proof here.
\bibitem{acquired} Human Compression project,
\emph{Residue stabilizers and distinguished generators};
\emph{Kernels and intersections in quotients of hypertopes}.
Complete acquired proofs with their distinct reviews and precise inputs.
\end{thebibliography}
\end{document}
